\section{Conclusion}

In this paper, we have evaluated the performance of \methodname across two extremes of the forecasting spectrum: short-term high-data numerical forecasting and long-term low-data event-based prediction. Our experiments demonstrate that LLMs are versatile forecasters capable of matching or surpassing human benchmarks in both domains. Specifically, we show that \methodname outperforms human expert superforecasters in probabilistic event prediction (Brier score 0.246 vs 0.287) and demonstrates a robust zero-shot capacity for processing high-volume numerical data where human manual forecasting is unfeasible.

Our findings suggest that LLMs can serve as a unified framework for diverse forecasting tasks, bridging the gap between statistical pattern recognition and abstract heuristic reasoning. Future work should focus on scaling these evaluations to larger datasets, implementing rigorous controls for temporal data leakage, and exploring the integration of specialized reprogramming frameworks like \timellm to enhance the precision of numerical tasks. Ultimately, the deployment of LLMs in forecasting holds the potential to significantly augment human decision-making in an increasingly complex and data-driven world.
